\section{研究对象、计算细节与仓库复现}\label{app:repro/protocol}
\subsection{主要符号与结果入口}
\begin{longtable}{@{}p{.20\textwidth}p{.73\textwidth}@{}}
\caption{主要符号及计算对象}\\\toprule 符号 & 含义\\\midrule\endfirsthead
\toprule 符号 & 含义\\\midrule\endhead
$R,P$ & 原始资料及文件、行、片段、位置索引。\\
$L^T,L^C,\bot$ & 任务要求、已展示贡献与证据不足的弃权。\\
$O^T,O^C,O^E,O^M$ & 两维候选、核验会话边和工具字段的可用指示。\\
$Y(1),Y(0),\tau$ & 无 AI 独立测量的潜在结果与目标因果效应。\\
$p_i,q,E_i$ & 学生候选层级分布、预设课程分布和核验相邻边。\\
$x_i,w,S_i$ & 五指标、非负和为一的权重、完整条件下的点分。\\
$\epsilon,b_i,\mathcal U_i,\mathcal W_j$ & 名义改标参数、整数改标条数、联合指标及权重集合。\\
\bottomrule
\end{longtable}

主结果为逐回合批次 t3。冻结清单位于
\begin{center}\small\path{runtime/research/full-turns-v1-20260926T143225Z/t3/frozen/snapshot-manifest.json}\end{center}
其 SHA-256 为\par\noindent\small\path{7e885ceecf12a180349369ed9ddfc962ad71f1ea270f68b037a86a396430631d}\normalsize\par
聚合分析入口为 \path{results/final-analysis/summary.json}，其中保存输入哈希、分母、方法条件、分析代码哈希和图表索引。论文所用聚合段落与表格在 \path{paper/final_analysis_results.tex}；图表由实际分析结果绘制。源文件、结果与最终 PDF 的哈希由交付清单绑定，代码版本使用仓库交付说明中的固定提交。

\subsection{来源、输出状态和统计分母}
两学期 2770 条问答行经范围筛选得到 1225 行，在其中恢复 3522 个学生提问回合，合格文本回合 3515。结构异常的 7 回合保留来源，不进入模型统计。DOCX 正文、图片及转写另有来源清单，未经验证的跨文件身份与会话连接不参与本研究的统一计数。学生—学期表保留 401 个单元；只有具有任务候选的 356 个单元进入候选指标敏感性，其他 45 个单元保留无候选状态。

基础逻辑任务为 $3515\times2=7030$，有效输出 6929，失败 101；有效包括明确弃权。随机复核为 352 回合、704 个逻辑任务，其中有效 692；额外风险复核为 827 回合、1654 个逻辑任务，其中有效 1638。12 个合成控制产生 48 个有效基础/复核输出；两快模型各 40 个重复任务，共 80 个，其中有效 78。真实 API 尝试、重试、逻辑任务和真实学生样本分别计数。相同输入仅在实验、模型、阶段、提示和参数完全相同时复用校验成功的响应；重复性检查不复用首次响应。

按维度先形成双模型状态：所需输出缺失记技术失败；两标签（含空值）不同记分歧；共同空值记弃权；任务维度共同非空但任一输出标为不确定时保留不确定；其余相同非空标签形成候选。已复核回合优先采用两强模型的候选或共同弃权；只有强复核技术失败、而双快已有候选或共同弃权时，回退双快结果。强模型之间的分歧继续保留。该规则逐条与封存最终状态核对，相同随机集流程比较使用同一套规则。

两模型六级一致性只用双方非空记录，三阶将 L1--2、L3--4、L5--6 合并。学生整簇重抽样次数为 2000，保留同一学生内部依赖；跨学生共享缓存造成的簇间依赖未在这些区间中额外校正。同模型重复、共享缓存与同题辅助复评均保留记录身份和来源。随机层与风险层分别列示，风险层用于定位难例，不能按其比例推断全体错误率。

\subsection{敏感性计算的定义域}
全回合未知标签范围固定所有最终候选，把其他回合分别赋予允许的最低与最高值。若总回合为 $N$、候选数为 $n$、候选层级和为 $s$、候选高阶数为 $h$，则
\begin{equation}
\mathrm{ABL}_{\rm raw}\in\left[\frac{s+N-n}{N},\frac{s+6(N-n)}{N}\right],\qquad
\HOT\in\left[\frac hN,\frac{h+N-n}{N}\right].
\end{equation}
该构造假设未知回合具有潜在六级值，评价的是这一假设下的敏感性；它没有把弃权重新认定为已观测标签。春减秋范围采用春季下界减秋季上界及春季上界减秋季下界。

学生分数范围的对象是具有候选的学生—学期单元。允许改标条数取 $b_i=\lceil\epsilon n_i\rceil$，取 $\epsilon=0,.05,.10,.20$。实现用 ABL/HOT/DHI 各自变化不超过 $b_i/n_i$ 的区间，截断至 $[0,1]$；缺失 CTQ/MAB 取 $[0,1]$。这些边际区间的笛卡尔积是式\eqref{eq:jointset}的保守外包，允许部分不能同时出现的指标组合。权重乘子为 $1\pm .10$，枚举五项乘子的 32 个顶点后归一；对外包矩形在每个顶点求线性极值，得到外包范围。区间不重叠是稳定区分的充分条件；重叠可能来自实际不确定性，也可能来自外包的保守性。

同学期比较共 $\binom{257}{2}+\binom{99}{2}=37747$ 对。严格大于对方上界并超过 $10^{-10}$ 数值容差时计为可稳定区分。该计数不跨学期连接身份，也不形成真实能力排序。$\epsilon=.05$ 时，向上取整后的改标比例中位数为 50\%，最大为 100\%；这反映多数单元候选很少。

\subsection{预实验和合成计算对照}\label{app:simulation}
首问预实验的单位为 384 条文本，与正文逐回合总体不同。该预实验的双快共同非空任务配对为 210 对、157 个学生簇，六级 $\kappa=0.823$，学生整簇 95\% 区间为 $[0.750,0.889]$；强复核子集共同非空为 50 对、49 簇，$\kappa=0.646$，区间为 $[0.458,0.802]$。由于输入粒度、复核选择和模型条件有差异，这些结果保留为流程探索记录。其逐题输出与失败记录在私有版本目录中单列。

有限总体合成实验取 $N=400$、层规模 $[240,120,40]$、单位成本 $[1,2,3]$、总预算 150，重复 500 次。给定实验者已知的残差方差，整数动态规划得到分配 $[49,37,9]$。表\ref{app:cost}用于检查分配程序及方差公式，成本单位为合成单位。
\begin{table}[htbp]\centering\small
\caption{相同预算下的合成参考标签配置}\label{app:cost}
\begin{tabular}{@{}lccc@{}}\toprule
策略 & 分配 $[n_1,n_2,n_3]$ & RMSE & 名义95\%区间覆盖\\\midrule
等量 & [25,25,25] & 0.0380 & 93.6\%\\
按规模 & [60,30,10] & 0.0321 & 93.2\%\\
已知方差／成本最优 & [49,37,9] & 0.0285 & 95.8\%\\\bottomrule
\end{tabular}
\end{table}

另以 24 名合成学生、每人 8 条标签比较三套权重和平方根聚合，检查聚合选择是否改变分数。文本红队保留 7 个条件、3 个模型共 21 个输出，包含冗长、动词注入及复制答案等输入；该实验使用早期模型/提示条件，每条件每模型一次，不用于估计 t3 的攻击成功率。

\subsection{仓库入口与一条复现命令}
代码、聚合数据、图表及运行说明见仓库：
\begin{center}\small\url{https://github.com/HNNUYUXUAN/nanxing-aiv}\end{center}
复算时检出交付说明对应的固定版本，并核对聚合清单中的来源哈希。授权研究资料与公开聚合材料有不同复算范围。完成固定依赖安装后，在项目根目录运行
\begin{center}\small\texttt{.venv/Scripts/python.exe\ scripts/reproduce.py}\end{center}
即可用聚合结果重绘研究图表并运行确定性演示样例。\path{--authorized-inputs} 在已具备授权原始资料时重建逐回合聚合分析，\path{--notebooks} 进一步执行研究 Notebook。具体输入、输出和运行回执见 \texttt{docs/}复现说明\texttt{.md}。模型 API 批次重跑与上述确定性复算分离；复算无需再次付费调用模型。

Python 依赖由 \path{requirements-lock.txt} 固定。论文使用仓库已有的多文件 XeLaTeX/Biber 工具链排版：
\begingroup\small
\begin{verbatim}
powershell -NoProfile -ExecutionPolicy Bypass `
  -File scripts/build_latex.ps1 `
  -Source paper/main.tex -OutputDirectory build/paper
\end{verbatim}
\endgroup
正文与参考文献、附录分别计页。原始学生记录、研究标识、数据库及凭据仅保存在授权私有资料中；公开仓库提供聚合结果和合成示例。旧环境归档与研究包操作记录保存在内部文档中。
